% Source coverage: physical PDF pp. 30--35 (printed pp. 7--12), ending with Eq. (15.2.9).
% Equations: (15.2.1)--(15.2.9), plus intervening unnumbered displays.
% Footnotes: one (physical p. 32). Citations: Ref. 7. Uncertainties: none.

\subsection{Gauge Theory Lagrangians and Simple Lie Groups}
\label{sec:15.2}

The transformation rules of the gauge-field tensor $F^a_{\mu\nu}$ and
the matter fields $\psi$ and their gauge-covariant derivatives do not involve the
derivatives of the transformation parameters $\epsilon^a(x)$, so if the
Lagrangian is constructed solely from these ingredients, and if it is invariant
under global transformations with $\epsilon^a$ constant, then it is invariant
under gauge transformations with general position-dependent $\epsilon^a(x)$.
We therefore assume that the Lagrangian satisfies these conditions: that is,
\begin{equation}
 \mathcal L=\mathcal L\bigl(\psi,D_\mu\psi,D_\nu D_\mu\psi,\ldots,
 F^a_{\mu\nu},D_\rho F^a_{\mu\nu},\ldots\bigr)
 \tag{15.2.1}\label{eq:15.2.1}
\end{equation}
with the invariance condition:
\begin{align}
 &\frac{\partial\mathcal L}{\partial\psi_\ell}
 \ii(t_a)_\ell{}^m\psi_m
 +\frac{\partial\mathcal L}{\partial(D_\mu\psi_\ell)}
 \ii(t_a)_\ell{}^m(D_\mu\psi_m)\notag\\
 &\quad
 +\frac{\partial\mathcal L}{\partial(D_\nu D_\mu\psi_\ell)}
 \ii(t_a)_\ell{}^m(D_\nu D_\mu\psi_m)+\cdots
 +\frac{\partial\mathcal L}{\partial F^b_{\mu\nu}}
 C^b{}_{ca}F^c_{\nu\mu}\notag\\
 &\quad
 +\frac{\partial\mathcal L}{\partial D_\rho F^b_{\mu\nu}}
 C^b{}_{ca}D_\rho F^c_{\nu\mu}+\cdots=0 .
 \tag{15.2.2}\label{eq:15.2.2}
\end{align}
On the other hand, the Lagrangian may not depend on the gauge field itself,
except insofar as it appears in $F^a_{\mu\nu}$ and in gauge-covariant
derivatives $D_\mu$. In particular, a mass term
$-\frac12m^2_{ab}A^a_\mu A^{b\mu}$ is ruled out.

We shall concentrate now on the terms in the Lagrangian that depend only on
$F^a_{\mu\nu}$. Just as in electrodynamics, for any massless particle of
unit spin the Lagrangian must contain a free-particle term quadratic in
$\partial_\mu A^a_\nu-\partial_\nu A^a_\mu$, and gauge invariance
then dictates that this free-particle term should appear as part of a term
quadratic in the field-strength tensor $F^a_{\mu\nu}$. Lorentz invariance
and parity conservation dictate its form as
\begin{equation}
 \mathcal L_A=-\frac12g_{ab}F^a_{\mu\nu}F^{b\mu\nu}
 \tag{15.2.3}\label{eq:15.2.3}
\end{equation}
with a constant matrix $g_{ab}$. If we do not assume parity (or CP or T)
conservation, then we may also include in the Lagrangian a term
\[
 \mathcal L'_A=-\frac12\theta_{ab}\epsilon^{\mu\nu\rho\sigma}
 F^a_{\mu\nu}F^b_{\rho\sigma}
\]
with another constant matrix $\theta_{ab}$. This term is actually a
derivative, and therefore does not affect the field equations or the Feynman
rules. Such a term would, however, have non-perturbative quantum mechanical
effects, to be discussed in Section 23.6.

Before going on to consider the properties of the matrix $g_{ab}$, it is
worth drawing attention to the fact that it is not possible to introduce a
kinematic term for the gauge field $A^a_\mu(x)$ without also including
interactions, the terms in Eq.~\eqref{eq:15.2.3} arising from the quadratic part
of the field strength $F^a_{\mu\nu}$ defined by
Eq.~\eqref{eq:15.1.13}. This is one more respect in which non-Abelian gauge
theories resemble general relativity, where the kinematic part of the Lagrangian
for the gravitational field is contained in the Einstein--Hilbert Lagrangian
density $-\sqrt{g}R/8\pi G$, which also contains self-interactions of the field.
The reasons in both cases are similar: the gravitational field interacts with
itself because it interacts with anything that carries energy and momentum, and
the gauge field interacts with itself because it interacts with anything that
transforms according to a non-trivial representation (in this case the adjoint
representation) of the gauge group. This is in contrast to the case of
electrodynamics, where the photon does not carry electric charge, the quantum
number with which it interacts, and it is consequently possible to introduce a
kinematic term $-\frac14F_{\mu\nu}F^{\mu\nu}$ for the electromagnetic field that
does not entail interactions.

The numerical matrix $g_{ab}$ may be taken symmetric, and must be taken
real to give a real Lagrangian. In order for this term to satisfy the
gauge-invariance requirement \eqref{eq:15.2.2}, we must have for all $u$:
\[
 g_{ab}F^a_{\mu\nu}C^b{}_{cd}
 F^{c\mu\nu}=0 .
\]
In order for this to be true without having to impose any functional relations
among the $F$s, the matrix $g_{ab}$ must satisfy the condition:
\begin{equation}
 g_{ab}C^b{}_{cd}
 =-g_{cb}C^b{}_{ad} .
 \tag{15.2.4}\label{eq:15.2.4}
\end{equation}
There is one more important condition on the matrix $g_{ab}$. Just as in
quantum electrodynamics, the rules of canonical quantization and the positivity
properties of the quantum mechanical scalar product require that the matrix
$g_{ab}$ in the Lagrangian \eqref{eq:15.2.3} must be positive-definite.
(That is, $g_{ab}u^a u^b$ is positive for all real $u$, and
vanishes for some real $u$ only if $u^a=0$ for all $a$.) This is
analogous to the requirement that in the kinematic Lagrangian
$-\frac12Z\partial_\mu\phi\partial^\mu\phi-\frac12m^2\phi^2$ for a real scalar
field $\phi$, the constant $Z$ must be positive-definite.

These requirements on the matrix $g_{ab}$ have far reaching
implications. They form one of a set of three equivalent conditions:

\begin{description}
 \item[\textbf{a:}] There exists a real symmetric positive-definite matrix
 $g_{ab}$ that satisfies the invariance condition
 \eqref{eq:15.2.4}.

 \item[\textbf{b:}] There is a basis for the Lie algebra (that is, a set of
 generators $\widetilde t_a=\mathcal S_{ab}t_b$, with
 $\mathcal S$ a real non-singular matrix) for which the structure constants
 $\widetilde C^a{}_{bc}$ are antisymmetric not only in the lower
 indices $b$ and $c$ but in all three indices $a$, $b$ and
 $c$. (In this basis it is convenient to drop the distinction between
 upper and lower indices $a$, $b$, etc., and write
 $\widetilde C_{abc}$ in place of
 $\widetilde C^a{}_{bc}$.)

 \item[\textbf{c:}] The Lie algebra is the direct sum of commuting compact
 simple and $U(1)$ subalgebras.\footnote{Some definitions: A subalgebra
 $\mathcal H$ of a Lie algebra $\mathcal G$ is a linear space, spanned by certain
 real linear combinations $t_i=\mathcal S_{ia}t_a$ of the generators
 $t_a$ of $\mathcal G$, such that $\mathcal H$ is itself a Lie algebra, in
 the sense that the commutators of the $t_i$ with each other are of the form
 $[t_i,t_j]=ic^k{}_{ij}t_k$. A subalgebra $\mathcal H$ is called invariant if
 the commutator of any element of the whole algebra $\mathcal G$ with any
 element of the subalgebra $\mathcal H$ is in the subalgebra $\mathcal H$. A
 simple Lie algebra is one without invariant subalgebras. A $U(1)$ subalgebra
 of $\mathcal G$ is one with just a single generator that commutes with all
 generators of the whole algebra $\mathcal G$. A semi-simple Lie algebra is one
 that has no invariant Abelian subalgebras, i.e., invariant subalgebras whose
 generators all commute with each other. Semi-simple Lie algebras are direct
 sums of simple (but not $U(1)$) Lie algebras. A simple or semi-simple Lie
 algebra is said to be compact if the matrix
 $\operatorname{Tr}\{t^A_a t^A_b\}
 =-C^c{}_{ad}C^d{}_{bc}$ is positive-definite.
 The meaning and importance of the properties of simplicity and compactness
 will be discussed further below. In saying that a Lie algebra $\mathcal G$ is
 a direct sum of subgroups $\mathcal H_n$, it is meant that it is possible to
 find a basis for $\mathcal G$ with generators $t_{na}$ for which the structure
 constants take the form
 \[
 C^{\ell c}{}_{na\,mb}=\delta_{\ell m}\delta_{mn}C^{(n)c}{}_{ab} ,
 \]
 where $C^{(n)c}{}_{ab}$ is the structure constant of the subalgebra
 $\mathcal H_n$.}
\end{description}

Appendix A of this chapter presents a proof of the equivalence of the conditions
\textbf{a}, \textbf{b}, and \textbf{c}~\hyperref[ch15-ref-7]{[7]}.

Before going on to discuss the physical implications of this result, it will be
useful to say a bit more about the condition of compactness. We will not use this
here, but a compact Lie algebra consists of the generators of a compact Lie group:
one for which the invariant volume of the group is finite. For instance, the
rotation group is compact; the Lorentz group is not. As a simple example of a
simple Lie algebra that is not compact, consider the commutation relations
\[
 [t_1,t_2]=-\ii t_3,\qquad [t_2,t_3]=\ii t_1,\qquad [t_3,t_1]=\ii t_2 .
\]
The structure constant here is real, but not completely antisymmetric; its
non-vanishing components are
\[
 C^3{}_{12}=-C^3{}_{21}=-1,\qquad
 C^1{}_{23}=-C^1{}_{32}=1,\qquad
 C^2{}_{31}=-C^2{}_{13}=1 .
\]
The metric given by Eq. (15.A.10) is here diagonal, with elements:
\[
 g_{11}=g_{22}=-g_{33}=-2 .
\]
This is not a positive matrix, so the Lie algebra is not compact. It is in fact
the Lie algebra of the non-compact group $O(2,1)$, the Lorentz group in two space
and one time dimensions.

Two sets of generators that differ by a real non-singular linear transformation
are considered to span the same Lie algebra, and generate the same group. This is
not true for complex linear transformations of generators. In particular, any
simple Lie algebra can be put into a compact form by a change of phase of the
generators in a suitable basis. For instance, for the Lie algebra of the above
example, it is only necessary to define new generators $t'_1=\ii t_1$, $t'_2=\ii t_2$,
$t'_3=t_3$, for which the commutation relations are
\[
 [t'_1,t'_2]=\ii t'_3,\qquad [t'_2,t'_3]=\ii t'_1,\qquad [t'_3,t'_1]=\ii t'_2 .
\]
The structure constant is now real and totally antisymmetric:
$C^a{}_{bc}=\epsilon_{abc}$. Here $g_{ab}=2\delta_{ab}$, and the algebra is
compact. We recognize this, of course, as the familiar algebra of the compact
group $O(3)$ of rotations in three dimensions. To see that this is always
possible for any simple Lie algebra, note that the matrix $g_{ab}$ defined by
Eq. (15.A.10) is real, symmetric, and non-singular, so that by a real orthogonal
transformation it may be put in a diagonal form with non-zero elements along the
main diagonal. It is then only necessary to multiply all the generators that
correspond in this basis to the negative diagonal elements of $g_{ab}$ by factors
$i$.

We note without proof that the finite-dimensional representations of compact Lie
groups are all unitary, and the finite-dimensional representations of compact Lie
algebras are correspondingly all Hermitian. Furthermore, it is easy to see that
the only Lie algebras that can have any non-trivial representation by independent
finite-dimensional Hermitian matrices $t_a$ are direct sums of $U(1)$ and
compact simple Lie algebras. To show this, we may simply define
\[
 g_{ab}=\operatorname{Tr}\{t_a t_b\} .
\]
This matrix is obviously positive-definite, because
$g_{ab}u^a u^b=\operatorname{Tr}\{(u^a t_a)^2\}$
is positive for any real $u^a$ and vanishes only if $u^a t_a=0$,
which is not possible unless all $u^a$ vanish because the $t^a$ are
assumed independent. Furthermore this $g_{ab}$ satisfies the invariance
condition \eqref{eq:15.2.4}, as can be seen by multiplying the commutation
relation \eqref{eq:15.1.2} with $t_d$ and taking the trace; this gives
\[
 iC^c{}_{ab}\operatorname{Tr}\{t_c t_d\}
 =\operatorname{Tr}\{[t_a,t_b]t_d\}
 =\operatorname{Tr}\{t_d t_a t_b-t_b t_a t_d\} ,
\]
which is obviously antisymmetric in $b$ and $d$. Having verified
\textbf{a}, we can rely on the above theorem to infer condition \textbf{c}, so
that the Lie algebra must be a direct sum of compact simple and $U(1)$
subalgebras.

Let's now return to the physics of gauge theories. In this section we have
inferred the existence of a positive symmetric real matrix $g_{ab}$,
that satisfies the invariance condition \eqref{eq:15.2.4}, from the necessity of
constructing a suitable kinematic term in the Lagrangian for the gauge field, and
in Appendix A of this chapter we have shown that this result is equivalent to a
condition on the Lie algebra, that it is a direct sum of compact simple and
$U(1)$ subalgebras. For our purposes, the important thing about this result is
that the simple Lie algebras are all of certain limited types and dimensionalities.
For instance, it is easy to see that these is no simple Lie algebra with less than
three generators, because in one or two dimensions there can be no non-zero
totally antisymmetric structure constants with three indices. With three
generators, an invariant subalgebra can be avoided by taking $C^3{}_{12}$,
$C^2{}_{31}$, and $C^1{}_{23}$ all non-zero. In the basis in which the structure
constant is real and totally antisymmetric, there is obviously only one
possibility:
\[
 C_{abc}=c\epsilon_{abc} .
\]
Here $c$ is an arbitrary non-zero real constant, which can be eliminated by a
change of scale of the generators, $t_a\longrightarrow t_a/c$, so the
Lie algebra is
\[
 [t_a,t_b]=\ii\epsilon_{abc}t_c .
\]
This may be recognized as the Lie algebra of the three-dimensional rotation group
$O(3)$, and also of the group $SU(2)$ of unitary unimodular matrices in two
dimensions, and was used as a basis for the original non-Abelian gauge theory of
Yang and Mills. Continuing in the same way, it can be shown that there are no
simple Lie algebras with 4, 5, 6, or 7 generators, one with 8 generators, and so
on. Mathematicians (notably Killing and E. Cartan) have been able to catalog all
simple Lie algebras. The compact forms of the simple Lie algebras form several
infinite classes of algebras of the `classical' Lie groups---the unitary
unimodular, unitary orthogonal, and unitary symplectic groups---plus just five
exceptional Lie algebras. This catalog is presented in Appendix B of this chapter.

It is also shown in Appendix A that under the equivalent conditions \textbf{a},
\textbf{b}, or \textbf{c}, the metric takes the form
\begin{equation}
 g_{ma,nb}=g_m^{-2}\delta_{mn}\delta_{ab}
 \tag{15.2.5}\label{eq:15.2.5}
\end{equation}
with real $g_m$, where $m$ and $n$ label the simple or $U(1)$ subalgebras, and
$a$ and $b$ label the individual generators of these subalgebras. We can
eliminate the constants $g_m^{-2}$ by a rescaling of the gauge fields
\begin{equation}
 A^\mu_{ma}\longrightarrow\widetilde A^\mu_{ma}
 \equiv g_m^{-1}A^\mu_{ma} ,
 \tag{15.2.6}\label{eq:15.2.6}
\end{equation}
but then in order to keep the same formulas \eqref{eq:15.1.10} and
\eqref{eq:15.1.13} for $D_\mu\psi$ and $F_a^{\mu\nu}$, we must also
redefine the matrices $t_a$ and the structure constants
\begin{equation}
 t_{ma}\longrightarrow\widetilde t_{ma}\equiv g_m t_{ma} ,
 \tag{15.2.7}\label{eq:15.2.7}
\end{equation}
\begin{equation}
 C^{(m)c}{}_{ab}\longrightarrow\widetilde C^{(m)c}{}_{ab}
 \equiv g_m C^{(m)c}{}_{ab} .
 \tag{15.2.8}\label{eq:15.2.8}
\end{equation}
That is, we can always define the scale of the gauge fields (now dropping the
tildes) so that $g_m$ in Eq.~\eqref{eq:15.2.5} is unity:
\begin{equation}
 g_{ab}=\delta_{ab} ,
 \tag{15.2.9}\label{eq:15.2.9}
\end{equation}
but then the transformation matrices $t_a$ and the structure constants
$C_{abc}$ contain an unknown multiplicative factor $g_m$ for each
simple or $U(1)$ subalgebra. These factors are the coupling constants of the gauge
theory. Alternatively, it is sometimes more convenient to adopt some fixed though
arbitrary normalization for the $t_a$ and structure constants within each
simple or $U(1)$ subalgebra, in which case the coupling constants appear in the
gauge-field Lagrangian \eqref{eq:15.2.3} as the factors $g_m^{-2}$ in
Eq.~\eqref{eq:15.2.5}.
